\cleardoublepage
\phantomsection
\addcontentsline{toc}{chapter}{ABSTRAK}
\begin{center}
{\FontJudulHalaman\bfseries ABSTRAK\par}
\end{center}
\vspace{12pt}
\begin{spacing}{1.5}
Kerja praktik di Divisi Maintenance PT XYZ menghasilkan sistem monitoring mesin cutting dan panel listrik berbasis Industrial Internet of Things (IIoT). Pada mesin cutting, Raspberry Pi membaca suhu motor, arus tiga fasa, tekanan udara, dan aliran udara melalui RS485/Modbus RTU. Hasil pengukuran ditampilkan pada HMI, disimpan dalam basis data SQLite, dan dikirim ke dashboard melalui MQTT. Pada monitoring panel, satu Raspberry Pi melayani empat panel di Power Room melalui modul MAX3485 pada bus RS485. PM2120 menyediakan data kelistrikan, sedangkan THM30MD menyediakan data suhu dan kelembaban. Program panel mengolah data pengukuran, memeriksa batas alarm, serta mengelola pembaca kartu RFID Wiegand yang dilengkapi buzzer dan LED terintegrasi. Dashboard panel menggunakan Django untuk menyediakan tampilan parameter, riwayat pengukuran, alarm, statistik, ekspor data, dan pengaturan. Sistem cutting dan panel mempunyai aplikasi serta penyimpanan data masing-masing. Penerapan satu Raspberry Pi untuk empat panel di Power Room menjadi dasar rancangan pengembangan monitoring per blok panel.
\end{spacing}
\vspace{6pt}
\noindent\textbf{Kata kunci:} IIoT, mesin cutting, panel listrik, Raspberry Pi, Modbus RTU, MQTT.
